\documentclass[12pt]{article}
\usepackage[T1]{fontenc}
\usepackage{lmodern}
\input{glyphtounicode}
\pdfgentounicode=1
\usepackage{amsmath, amsthm, amssymb}
\usepackage[margin=1in]{geometry}
\usepackage{hyperref}
\usepackage{mathtools}
\usepackage{doi}
\hypersetup{
  pdftitle={What Weight Stabilisation Would Require: A Finite-Window Audit of Admissibility Weights},
  pdfauthor={J\'er\^ome Beau},
  pdfsubject={Conditional weight convergence and finite-window calibration},
  pdfkeywords={admissibility weights, finite BFS window, fingerprint response, withdrawn claim}
}

\providecommand{\doi}[1]{\href{https://doi.org/#1}{\texttt{doi:#1}}}

\newtheorem{theorem}{Theorem}
\newtheorem{lemma}[theorem]{Lemma}
\newtheorem{proposition}[theorem]{Proposition}
\newtheorem{corollary}[theorem]{Corollary}
\theoremstyle{definition}
\newtheorem{definition}[theorem]{Definition}
\newtheorem{hypothesis}[theorem]{Hypothesis}
\theoremstyle{remark}
\newtheorem{remark}[theorem]{Remark}

\DeclareMathOperator{\Heis}{Heis}
\DeclareMathOperator{\op}{op}
\newcommand{\cBI}{c_{\mathrm{BI}}}
\newcommand{\sstar}{\sigma^{*}}
\newcommand{\Eq}{\mathcal{E}_q}
\newcommand{\Eoo}{\mathcal{E}}
\newcommand{\nstar}{n^{*}(q)}

\title{What Weight Stabilisation Would Require:\\
A Finite-Window Audit of Admissibility Weights}
\author{J\'er\^ome Beau}
\date{2026\\[2pt]\small Version 2.0, candidate}

\begin{document}

  \maketitle

  \begin{abstract}
    Version 1 claimed to prove common positive convergence of the admissibility
    weights from spectral universality.
    That proof imported a relative profile estimate never supplied by U1, identified
    the weights with shell-capacity sums without a bridge to Q5a's fingerprint-response
    averages, and sent the calibrated fitting-window endpoint to infinity although
    the established five-block protocol has endpoint at most~22 for large primes.
    We retain the exact conditional summation inequality, but it concerns a separately
    defined partial-sum proxy.
    For a fixed nonnegative profile and the published window protocol, that proxy's
    target is the sum through depth~22, provided the necessary convergence input is
    supplied; it is not the infinite series asserted in version~1.
    Neither the proxy inequality nor fibre constancy proves convergence of Q5a's
    actual weights, whose horizontal-generator limits may differ.
    The finite O25 pair exponents do not establish a global pointwise profile bound
    or identify a Born--Infeld normalisation.
    Thus [H-w] remains open on these inputs.
    \textbf{Interpretive outlook.}
    The audit isolates a testable bridge: construct the measured response weights
    from the block-sampled rank observable and control them uniformly on the
    calibrated finite window before assigning an effective kinetic coefficient.
  \end{abstract}

  \paragraph{Keywords.}
    Spectral admissibility; Dirichlet form; finite BFS window; fingerprint response;
    Heisenberg group; conditional weight convergence; rank increment;
    withdrawn claim; Born--Infeld normalisation.

    \clearpage
    \tableofcontents

    \clearpage

  \section{Introduction}
  \label{sec:intro}

  The published W1 claim was that the four generator weights of the Q5a
  admissibility form converge to one positive constant.
  Its proof used a relative universality estimate attributed to U1.
  U1 version 2.0~\cite{Beau2026u1} withdraws even the absolute estimate claimed
  in its version 1; the relative estimate used here was never proved there.
  Moreover, Q5a version 3.2~\cite{Beau2026q5a} now states a different weight
  observable and allows distinct limits for the two generator directions.
  This paper records the surviving elementary implication and the missing
  identifications rather than asserting a closed chain.

  The numerical source must be read at its own scope.
  O25~\cite{Beau2026a29} fits pair exponents over finite windows.
  The historical value $\delta_{\mathrm{pair}}\approx7.44$ came from O16's
  single-pair sample; O25's five-prime means range from 8.784 to 9.983, and its
  large-prime raw statistic reaches 7.61 at $q=601$.
  O25 expressly does not determine an asymptotic exponent.
  None of these values is a bound on every term of a limiting profile.

  \subsection*{Standing conventions}

  Let $q$ range over odd primes, $G_q=\Heis_3(\mathbb Z/q\mathbb Z)$,
  and $S_q=\{X^{\pm1},Y^{\pm1}\}$.
  Let $n^*(q)$ denote the upper endpoint of the calibrated O25/O28 fitting
  window when that protocol is intended.
  The O25 observable $\sigma_{q,\mathbf c}(n)=\Delta r_{q,\mathbf c}(n)/|S_n|$
  depends on a three-component generic block $\mathbf c$ and a sampling rule,
  not only on one central character~\cite{Beau2026a29}.
  We use $p(n)$ for an arbitrary supplied nonnegative comparison profile.
  No theorem below identifies $p$ with a limit of the O25 observable.
  In Q5a, $a_q(s)$ denotes the actual fingerprint-response average.
  A partial-sum proxy will be written $b_q(s)$ to keep the two objects distinct.
  \section{Admissibility Weights and Hypothesis [H-w]}

  Q5a version 3.2~\cite{Beau2026q5a} calls $a_q(s)$ a fingerprint-response
  average and uses the form
  \[
    \Eq(f)=\frac{q^{-2}}{2}\sum_{s\in S_q}a_q(s)
    \|(\rho_q(s)-I)f\|^2
  \]
  on its filtration stage, without an interior projector.
  Its current [H-$w'$] asks separately for positive limits $A_X$ and $A_Y$
  for the $X$ and $Y$ generator pairs.
  A common limit is a stronger claim and requires its own evidence.

  \begin{definition}[Partial-sum proxy, distinct from Q5a weights]
    \label{def:aw}
    Suppose a family of nonnegative shell observables
    $\sigma_{q,c}(n)$ and a map $c_q:S_q\to(\mathbb Z/q\mathbb Z)^\times$
    are supplied.
    Define the \emph{proxy}
    \begin{equation}
      \label{eq:aw}
      b_q(s):=\sum_{n=1}^{n^*(q)}\sigma_{q,c_q(s)}(n).
    \end{equation}
    Neither the map $c_q$ nor an identity $a_q(s)=b_q(s)$ is supplied by
    the O25 or Q5a statements cited here.
    In particular, the displayed definition is not Q5a Definition~2.3.
  \end{definition}

  \begin{remark}[Partial sum versus Ces\`aro mean]
    \label{rem:partial-sum}
    For a fixed nonnegative summable profile $p$ with
    $A_\infty:=\sum_{n\ge1}p(n)>0$ and an unbounded cutoff $N$, the mean
    $N^{-1}\sum_{n\le N}p(n)\sim A_\infty/N$, while the partial sum tends
    to $A_\infty$.
    The former $O(N^{1-\delta^*/2})$ assertion for a profile
    $p(n)\asymp n^{-\delta^*/2}$ with exponent greater than~2 was false.
    For the calibrated finite window, neither calculation identifies the
    actual Q5a weight or its Mosco limit.
  \end{remark}

  \begin{hypothesis}[Historical common-weight hypothesis [H-w]]
    \label{hyp:Hw}
    The version-1 claim required a constant $A>0$ such that
    \begin{equation}
      \label{eq:Hw}
      \sup_{s\in S_q}|a_q(s)-A|\longrightarrow0.
    \end{equation}
    This is retained as an open target, not as a result of this paper.
    It implies the pair of Q5a [H-$w'$] limits with $A_X=A_Y=A$,
    but Q5a's separate limits do not imply a common value.
  \end{hypothesis}

  \section{Main Result}

  \begin{theorem}[Conditional proximity of the partial-sum proxy]
    \label{thm:proximity}
    Suppose Definition~\ref{def:aw} applies and a supplied profile
    $p(n)\ge0$ satisfies, for an error $\varepsilon_q\ge0$,
    \begin{equation}
      \label{eq:U1}
      |\sigma_{q,c}(n)-p(n)|\le\varepsilon_q p(n)
      \quad\text{for every relevant }c,\quad 1\le n\le n^*(q).
    \end{equation}
    Then, with $A_q:=\sum_{n=1}^{n^*(q)}p(n)$,
    \begin{equation}
      \label{eq:proximity}
      \sup_{s\in S_q}|b_q(s)-A_q|\le\varepsilon_q A_q.
    \end{equation}
    U1 version 2.0 supplies neither~\eqref{eq:U1} nor a value
    $\varepsilon_q=O(q^{-1/2})$.
  \end{theorem}

  \begin{proof}
    For each $s$, sum the absolute deviations in~\eqref{eq:U1} over
    $1\le n\le n^*(q)$ and apply the triangle inequality.
    The same bound holds after taking the supremum over $S_q$.
    The proof uses no property of Q5a's response averages.
  \end{proof}

  \begin{proposition}[Fibre constancy does not give common weights]
    \label{prop:fibre}
    Constancy of an observable on each fibre of a non-injective
    projection does not imply that its values are equal across fibres
    or that they converge uniformly as a parameter varies.
  \end{proposition}

  \begin{proof}
    Let $\Omega=\{a,a',b,b'\}$ and let $\Pi$ map $a,a'$ to~0 and
    $b,b'$ to~1.
    Set the observable equal to~0 on the first fibre and~1 on the
    second for every $q$.
    It is fibre-constant, yet its inter-fibre spread stays equal to~1.
    Conjugate-pair parity, even if established for the measured
    observable, equates at most the two members of each pair and
    does not supply equality across distinct pairs.
  \end{proof}

  O17's exact block-independence theorem concerns its explicit toy
  model~\cite{Beau2026a21}; it is not an identity for the independently
  sampled three-component O25 rank observable.
  O18 records the parity-to-fibre identification as open.
  The version-1 transfer of these results to all O25 characters and
  then to the four Q5a weights is therefore an unproved bridge.

  \begin{lemma}[Target of a bounded fitting window]
    \label{lem:Aq}
    Let $p(n)\ge0$ be a fixed profile and
    $A_q=\sum_{n=1}^{n^*(q)}p(n)$.
    If $n^*(q)\to22$ in probability, then
    \begin{equation}
      \label{eq:Aq-limit}
      A_q\longrightarrow A_{22}:=\sum_{n=1}^{22}p(n)
      \quad\text{in probability}.
    \end{equation}
    The limit is positive if at least one of its first 22 terms is
    positive.
    No summability beyond depth~22 is needed.
    Neither monotonicity of $A_q$ in $q$ nor equality with
    $\sum_{n\ge1}p(n)$ follows.
  \end{lemma}

  \begin{proof}
    Because $n^*(q)$ is integer-valued, convergence in probability
    to~22 gives $\mathbb P(n^*(q)=22)\to1$.
    On that event $A_q=A_{22}$ exactly.
    Positivity is immediate from the stated nonnegative term.
  \end{proof}

  The Critical Coverage note~\cite{Beau2026ccnote} proves
  $n_1(q)\le22$ for every prime $q\ge311$ and $n_1(q)\to22$ in
  probability for its stated five-block calibration.
  Version 1 identifies $n^*(q)$ with the O25/O28 calibrated endpoint.
  Applying the lemma to another sampling protocol requires identifying
  its endpoint with that five-block statistic; it cannot be assumed.
  In particular, the old claim that the calibrated endpoint grows at a
  polynomial rate contradicts the proved five-block law.

  \begin{corollary}[Conditional bridge to a common weight]
    \label{cor:Hw}
    Suppose the hypotheses of Theorem~\ref{thm:proximity} hold with
    $\varepsilon_q\to0$, the same cutoff obeys $n^*(q)\to22$ in
    probability, and the actual Q5a weights satisfy
    $d_q:=\sup_{s\in S_q}|a_q(s)-b_q(s)|\to0$ in probability.
    Then
    \begin{equation}
      \label{eq:Hw-rate}
      \sup_{s\in S_q}|a_q(s)-A_{22}|
      \le d_q+\varepsilon_q A_q+|A_q-A_{22}|
      \longrightarrow0
      \quad\text{in probability}.
    \end{equation}
    If the three premises hold deterministically with an eventually
    constant cutoff, the conclusion holds deterministically.
    For $A_{22}>0$, that deterministic conclusion would imply
    Hypothesis~\ref{hyp:Hw}.
    None of the three load-bearing premises is proved here for Q5a.
  \end{corollary}

  \begin{proof}
    Insert $b_q(s)$ and $A_q$, use Theorem~\ref{thm:proximity},
    then apply Lemma~\ref{lem:Aq}.
  \end{proof}

  \begin{remark}[Status of the structural argument]
    \label{rem:PTO}
    Non-injectivity says that an observable on $\mathrm{Im}\,\Pi$ is
    constant on each projection fibre.
    It does not identify the O25 block observable with Q5a's generator
    response averages, prove cross-fibre uniformity, or select a common
    limit for $X$ and $Y$.
    The projective temporal-ordering analogy is interpretive and
    supplies none of these estimates.
  \end{remark}

  \section{Can the Limit Be Identified with Born--Infeld Saturation?}
  \label{sec:A}

  \subsection{The measured object and a proposed normalisation}

  The O25 quantity is a Gram--Schmidt rank increment divided by a shell
  cardinality, evaluated on sampled three-component blocks.
  A Born--Infeld bound on a projected operator norm does not, by itself,
  identify or bound that rank increment term by term.
  The compressed operator-norm profile formerly imported from U1 is not
  the O25 observable; U1 version 2.0 explicitly withdraws that
  identification~\cite{Beau2026u1}.
  Thus the claimed formula
  $\sigma^*(n)=\|\Pi_{S_n}d\pi_1\Pi_{S_n}\|_{\op}$ is not available.

  Bass growth concerns balls in the infinite discrete Heisenberg group.
  A statement about finite $G_q$ shells needs a specified pre-wrap range;
  it cannot be asserted as a global $n\to\infty$ asymptotic for a finite
  group.
  Even if a shell law proportional to $n^3$ applies in such a range,
  its reciprocal gives an $n^{-3}$ leading profile only after a separate
  saturation and norm-to-rank bridge.
  It cannot simultaneously have a leading $n^{-p}$ law with
  $p\approx4.75$--$5$: subleading corrections do not change the
  leading exponent.

  \subsection{What remains of the explicit formula}

  The former expression
  \begin{equation}
    \label{eq:A-formula}
    A_\infty\approx
    \frac{c_{\mathrm{BI}}^2}{C_{\Heis}}\,
    \zeta\!\left(\frac{\delta^*}{2}\right)
  \end{equation}
  is retained only as the version-1 proposal under a supplied global
  power-law profile, saturation and normalisation.
  It is not a consequence of O25, Q5a or U1.
  For a bounded calibrated window the candidate target would instead be
  $A_{22}=\sum_{n=1}^{22}p(n)$, if the proxy assumptions of
  Corollary~\ref{cor:Hw} could be established.
  No numerical value or Born--Infeld identity for Q5a's actual weights
  follows from that conditional expression.

  \begin{remark}[Status of the proposed formula]
    \label{rem:formula}
    Equation~\eqref{eq:A-formula} is an unverified model expression, not
    an identification of an admissibility coefficient.
    To turn it into one requires a construction of the Q5a weight from
    the O25 observable, a proved profile with controlled normalisation,
    a valid shell range, and an independent saturation principle.
    Fitting a finite pair exponent supplies none of these bridges.
  \end{remark}

  \section{Consequences for the Q5a Programme}
  \label{sec:consequences}

  \subsection{Status of the weight claim}

  The historical Q5a formulation used the common-weight condition
  [H-w] in a proposed Mosco argument.
  Current Q5a version 3.2 defines the response-average weights and
  states separate positive limits [H-$w'$] for the $X$ and $Y$ pairs.
  It also identifies a zero-form limit at the published normalisation
  under its stated conditions.
  W1 proves neither the historical common-weight assertion nor the
  current pair of limits.

  \begin{table}[ht]
    \centering
    \begin{tabular}{p{0.23\textwidth}p{0.62\textwidth}}
      \hline
      Claim & Status after this audit \\
      \hline
      Proxy bound & Theorem~\ref{thm:proximity}, conditional on a supplied
        relative profile estimate and a generator-to-character map. \\
      Window target & Lemma~\ref{lem:Aq}, conditional on using the same
        bounded five-block endpoint and a supplied profile. \\
      Q5a weights & The proxy-to-response identification is open;
        [H-$w'$] and the stronger common [H-w] remain open. \\
      Born--Infeld value & No coefficient is identified from the present
        bound or the finite O25 fits. \\
      \hline
    \end{tabular}
    \caption{The proxy statements do not close a Q5a weight hypothesis.}
    \label{tab:status}
  \end{table}

  \subsection{Dependence of later papers}

  H2, Q5a-O2 and the Emergent Geometry Note cite the former W1 proof
  when classifying [H-w] as established.
  Those uses must be retyped during their own release steps.
  Q8, Q10 and Q11 may retain [H-w] as an explicit hypothesis, but may
  not discharge it by citing W1.
  No conclusion about the Q10 spatial coefficient or the Q11 temporal
  coefficient follows from the present proxy calculation.

  \section{Finite O25 Data and Their Scope}

  The five-prime O25 campaign~\cite{Beau2026a29} reports the following
  means over conjugate pairs, with inter-pair standard deviations:
  \begin{center}
    \begin{tabular}{ccc}
      \hline
      $q$ & fitted $\delta_{\mathrm{pair}}$ mean & inter-pair s.d. \\
      \hline
      29  & 8.893 & 0.539 \\
      61  & 9.983 & 0.259 \\
      101 & 9.548 & 0.161 \\
      151 & 9.125 & 0.139 \\
      211 & 8.784 & 0.109 \\
      \hline
    \end{tabular}
  \end{center}
  O25 extends the campaign to $q=307,401,601$ with fewer block samples
  per pair and reports a raw global exponent of 7.61 at $q=601$.
  The historical 7.44 value is an O16 single-pair estimate, not a
  profile bound or the determined O25 asymptote.
  These measurements show finite-$q$ pair concentration under their
  sampling protocol.
  They do not estimate $\sup_s|a_q(s)-b_q(s)|$, establish a relative
  $q^{-1/2}$ profile rate, or prove that a nonnegative infinite profile
  is summable.

  \section{Conclusion}

  Version 1's proof of a common positive admissibility weight is
  withdrawn.
  The summation step remains exact when its relative profile premise
  and proxy definition are explicitly supplied.
  With the published five-block calibration, the natural partial-sum
  target is a finite-depth quantity rather than the former infinite
  series, but even that observation concerns a proxy, not Q5a's
  fingerprint-response averages.

  \paragraph{Interpretive outlook.}
  A common kinetic coefficient would express a measurable agreement of
  the two generator directions.
  The current evidence does not select it.
  The productive next test is to construct the response weights from
  the block-sampled rank data and check their directional limits on the
  actual calibrated window, keeping any Born--Infeld normalisation as
  a separately tested identification.

  \section*{Acknowledgements}

  The author thanks the Cosmochrony programme for the O25 numerical data.

  \bibliographystyle{plain}
  \bibliography{cosmochrony-bibliography,external-refs}

\end{document}
